%% kapa: exam=1 round=24 year=2013 period=1 form=A subject=economics subject_ko=경제원론 seq=30 number=70
%% source: Q-Net 감정평가사 기출문제 (cst00309) artlSeq=5073802 file=extracted/1008/5073802/2013년 제24회 감정평가사 1차시험 기출문제/1교시(A형).hwp
%% conversion: kordoc 4.24.51 -> tools/gen1.py (auto); review=proofread-sonnet; flags=none
\begin{question}{70}
구매력평가설이 성립할 때, 다음 설명 중 옳지 않은 것은?
\begin{choices}
  \item 자국의 통화량이 증가할 때 실질환율은 변화하지 않는다.
  \item 외국의 양적완화 정책으로 외국의 물가가 상승하면, 자국의 실질 순수출이 증가한다.
  \item 양국 물가상승률의 차이가 명목환율의 변화율에 영향을 준다.
  \item 양국간 무역에서 재정거래(arbitrage)에 의한 수익을 얻을 수 없다.
  \item 양국 물가수준의 상대적 비율이 명목환율에 영향을 준다.
\end{choices}
\end{question}
\answer{2}
